\makeatletter
\@ifundefined{keepfig}{\newenvironment{keepfig}{\par\edef\keepfig@pi{\the\parindent}\noindent
  \begin{minipage}{\linewidth}\parindent=\keepfig@pi\relax}{\end{minipage}\par}}{}
\makeatother
\section{Разрешимость уравнений}

\begin{definition}[3.27i]
Пусть $f\in F[X]$. Говорят, что уравнение $f(X)=0$ (или многочлен $f$) \emph{разрешимо в радикалах}, если его корни
получаются из элементов $F$ сложением, вычитанием, умножением, делением и извлечением корней $m$-й степени, точнее~---
если существует башня полей
\[
F=F_0\subset F_1\subset F_2\subset\dots\subset F_m
\]
такая, что
\begin{enumerate}[label=(\alph*)]
\item $F_i=F_{i-1}[\alpha_i]$, где $\alpha_i^{m_i}\in F_{i-1}$ для некоторого целого $m_i\ge1$;
\item $F_m$ содержит поле разложения $f$ (т.\,е.\ $f$ раскладывается в $F_m[X]$ на линейные множители).
\end{enumerate}
Башню со свойством (a) называют \emph{радикальной}.
\end{definition}

Многочлен $f$ здесь не обязан быть неприводимым: см.\ \rf{3.28b}. Условие (b) требует, чтобы поле разложения
\emph{содержалось} в $F_m$; совпадать с ним $F_m$ не обязано, и обычно не совпадает (\rf{3.28a}).

\begin{theorem}[3.28: Галуа, 1832]
Пусть $F$~--- поле характеристики $0$ и $f\in F[X]$. Уравнение $f(X)=0$ разрешимо в радикалах тогда и только тогда,
когда группа Галуа $G_f$ разрешима.
\end{theorem}

Доказательство дано в \rf{5.34}. Его идея: если $G_f$ разрешима, то после присоединения нужных корней из единицы
расширение $F_f/F$ разбивается по цепочке подгрупп с циклическими факторами на циклические шаги, а циклический шаг над
полем с корнями из единицы~--- присоединение радикала (теория Куммера, \rf{5.27}); обратно, нормальное замыкание
радикальной башни (с добавленными корнями из единицы) имеет разрешимую группу Галуа, и $G_f$~--- её факторгруппа.
Группа $S_n$ при $n\ge5$ неразрешима (знакопеременная группа $A_n$ проста и неабелева), поэтому многочлены с группой
Галуа $S_n$, $n\ge5$, не разрешимы в радикалах; такие многочлены над $\Q$ есть в любой степени $n$ (\rf{4.32}), например
$X^5-6X+3$ (\rf{5.35e}). Уравнения степени не выше $4$ разрешимы в радикалах (\rf{5.35d}).

\begin{example}[3.28a: формула Кардано]
Решим по Кардано уравнение $X^3+pX+q=0$, $p,q\in\CC$ (к такому виду любое кубическое уравнение $X^3+aX^2+bX+c=0$
приводится подстановкой $X=Y-\frac a3$). Ищем корень в виде $x=u+v$. Тогда $x^3=u^3+v^3+3uv(u+v)$, и уравнение принимает вид
$u^3+v^3+q+(3uv+p)(u+v)=0$; достаточно потребовать $3uv=-p$ и $u^3+v^3=-q$. Тогда $u^3,v^3$~--- корни уравнения
$T^2+qT-\frac{p^3}{27}=0$, и
\[
x=\sqrt[3]{-\frac q2+\sqrt{\frac{q^2}4+\frac{p^3}{27}}}+\sqrt[3]{-\frac q2-\sqrt{\frac{q^2}4+\frac{p^3}{27}}},\qquad
\frac{q^2}4+\frac{p^3}{27}=-\frac D{108},\quad D=-4p^3-27q^2,
\]
где $D$~--- дискриминант (решение упражнения 1-4), а кубические корни $u$, $v$ выбираются с произведением $-\frac p3$. Два других
корня~--- $\omega u+\omega^2v$ и $\omega^2u+\omega v$, $\omega=e^{2\pi i/3}$: при $1+\omega+\omega^2=0$ прямое перемножение даёт
\[
\bigl(X-(u+v)\bigr)\bigl(X-(\omega u+\omega^2v)\bigr)\bigl(X-(\omega^2u+\omega v)\bigr)=X^3-3uv\,X-(u^3+v^3)=X^3+pX+q .
\]

Для $f=X^3+3X+1\in\Q[X]$ (\rf{3.23c}: $f$ неприводим, $G_f\cong S_3$, $[\Q_f:\Q]=6$, единственное квадратичное подполе
$\Q_f$~--- это $\Q[\sqrt{-15}]$) имеем $p=3$, $q=1$, $D=-135$, $-\frac D{108}=\frac54$, $u^3,v^3=\frac{-1\pm\sqrt5}2$, и формула даёт
вещественные радикалы $u=\sqrt[3]{(-1+\sqrt5)/2}$, $v=-\frac1u$ из \rf{3.23c}, вещественный корень $x_1=u-\frac1u\approx-0{,}3222$ и
радикальную башню
\[
\Q\ \subset\ \Q[\sqrt5]\ \subset\ \Q[u]\ \subset\ \Q[u,\omega],\qquad (\sqrt5)^2=5,\quad u^3=\tfrac{-1+\sqrt5}2,\quad\omega^3=1,
\]
со степенями шагов $2$, $3$, $2$: поле $\Q[u]$ содержит $\sqrt5=2u^3+1$ (степени $2$) и $x_1$ (степени $3$), так что его степень
делится на $6$, а больше $6$ она быть не может, так как $u^3\in\Q[\sqrt5]$; а $\omega\notin\Q[u]\subset\R$. Поле $\Q[u,\omega]$
степени $12$ содержит все три корня, т.\,е.\ поле разложения $\Q_f$ степени $6$, но не совпадает с ним: башня <<выходит>> за
пределы $\Q_f$ уже на первом шаге, так как $\sqrt5\notin\Q_f$ (\rf{3.23c}). Можно начать и с корней из единицы, как в
доказательстве \rf{5.34}: $\Q\subset\Q[\omega]\subset\Q[\omega,\sqrt{-15}]\subset\Q[\omega,\sqrt{-15},u]$~--- сначала корень из
единицы, затем квадратный корень из дискриминанта, затем кубический радикал ($u^3\in\Q[\sqrt5]\subset\Q[\omega,\sqrt{-15}]$, так
как $\sqrt{-15}=\sqrt5\cdot\sqrt{-3}$ и $\sqrt{-3}=2\omega+1$); вершина та же, $\Q[u,\omega]$.

Здесь дискриминант отрицателен, и вещественный корень выражается вещественными радикалами. При трёх вещественных
корнях ($D>0$) так не получится: для $X^3-3X+1$ ($p=-3$, $q=1$, $D=81$) под квадратным корнем стоит $-\frac34$, и
формула даёт $x_1=2\cos\frac{2\pi}9=e^{2\pi i/9}+e^{-2\pi i/9}$ через кубические корни из невещественных чисел
$e^{\pm2\pi i/3}$ (\rf{3.22a}); по \rf{3.22f} это неизбежно (<<неприводимый случай>>).
\end{example}

\begin{remark}[3.28b]
(a) \emph{Неприводимость $f$ в определении \rf{3.27i} не нужна, и ничего не меняет.} В характеристике $0$ многочлен разрешим
в радикалах тогда и только тогда, когда разрешимы в радикалах все его неприводимые множители: по \rf{3.28} это
равносильно тому, что разрешимы группы Галуа, а $G_f$ разрешима тогда и только тогда, когда разрешимы группы Галуа всех
различных неприводимых множителей (\rf{3.27h}).

(b) \emph{Для неприводимого $f$ \textup(в характеристике $0$\textup) достаточно, чтобы в радикальной башне лежал один корень.} Пусть $\alpha\in F_m$~--- корень
неприводимого $f$. В доказательстве \rf{5.34} (часть <<$\Rightarrow$>>) по радикальной башне строится нормальное замыкание
$\tilde E\supset F_m$, которое само является вершиной радикальной башни над $F$; так как $\tilde E/F$ нормально и содержит
корень $\alpha$ неприводимого многочлена $f$, оно содержит все корни $f$ (\rf{3.10}), т.\,е.\ поле разложения $f$. Значит,
$f$ разрешим в радикалах.
\end{remark}

\subsection*{Характеристика $p$}

\begin{remark}[3.29]
Если $F$ имеет характеристику $p>0$, теорема \rf{3.28} неверна по двум причинам.
\begin{enumerate}[label=(\alph*)]
\item Многочлен $f$ может быть несепарабельным и не иметь группы Галуа в смысле \rf{3.27a}. Например, $X^p-t$ над $\F_p(t)$
неприводим (\rf{2.3}), а его производная равна нулю; его поле разложения $\F_p(t)[t^{1/p}]$ чисто несепарабельно,
хотя в смысле определения \rf{3.27i} он разрешим в радикалах.
\item Многочлен $X^p-X-a$ всегда сепарабелен (его производная равна $-1$; корни~--- $p$ различных элементов
$\alpha+i$, $i\in\F_p$, \rf{5.32c}(a)) и имеет абелеву группу Галуа~--- тривиальную или $\Z/p\Z$ (\rf{5.32c}(b)), но может
быть не разрешим в радикалах (\rf{3.29a}, \rf{5.32j}).
\end{enumerate}
Если изменить определение разрешимости, разрешив в башне шаги, задаваемые многочленами вида (b), и запретив корни
степени, делящейся на $p$ (корень степени $p$ даёт лишь чисто несепарабельный шаг, ср.\ (a)), то для сепарабельных
$f\in F[X]$ теорема остаётся верной в любой характеристике (\rf{5.35a}, \rf{5.35b}).
\end{remark}

Почему $X^p-X-a$ может быть не разрешим в радикалах? Если $a\notin\wp(F)$, где $\wp(x)=x^p-x$, то по \rf{5.32c} поле
$F[\alpha]$, $\alpha^p-\alpha=a$, есть расширение Галуа степени $p$. Наивное соображение <<тогда $\alpha$ должен быть
радикалом: $\alpha^m\in F$>> не годится: определение \rf{3.27i} допускает длинные башни с вспомогательными радикалами, и
требуется лишь, чтобы $F[\alpha]$ лежало где-то внутри $F_m$. Нужно найти величину, которую можно проследить вдоль всей
башни. Такая величина~--- сепарабельная степень (число вложений в алгебраически замкнутое поле, \rf{5.46a}).

\begin{proposition}[3.29a]
Пусть $\operatorname{char}F=p>0$, и $F$ содержит первообразный корень степени $q$ из единицы для каждого простого $q\ne p$
(например, $F=\Fpb(t)$). Если $a\in F\smallsetminus\wp(F)$, то ни один корень многочлена $X^p-X-a$ не лежит ни в какой
радикальной башне над $F$. В частности, многочлен $X^p-X-t$ с группой Галуа $\Z/p\Z$ над $\Fpb(t)$ не разрешим в радикалах.
\end{proposition}

\begin{proof}
Пусть $F=F_0\subset\dots\subset F_m$~--- радикальная башня, $F_i=F_{i-1}[\alpha_i]$, $\alpha_i^{m_i}\in F_{i-1}$, и
$\alpha\in F_m$, $\alpha^p-\alpha=a$.

\emph{Можно считать все $m_i$ простыми.} Шаг с $m_i=1$ тривиален. Если $m_i=q_1q_2\cdots q_s$ (простые с повторениями),
вставим между $F_{i-1}$ и $F_i$ поля $F_{i-1}[\alpha_i^{m_i/q_1}]\subset F_{i-1}[\alpha_i^{m_i/(q_1q_2)}]\subset\dots\subset F_{i-1}[\alpha_i]$;
на $j$-м из этих шагов присоединяется элемент, $q_j$-я степень которого лежит в предыдущем поле.

Пусть $\Om$~--- алгебраическое замыкание $F_m$ (глава 6); оно является и алгебраическим замыканием $F$. Для конечного
расширения $K/L$ подполей $\Om$ через $[K:L]_s$ обозначим сепарабельную степень~--- число продолжений вложения $L\to\Om$
до вложения $K\to\Om$; она не зависит от вложения, мультипликативна в башнях, для $K=L[\beta]$ равна числу различных
корней минимального многочлена $\beta$, и $[K:L]_s=[K:L]$ для сепарабельного $K/L$ (\rf{5.46a}).

\emph{Шаги башни.} Пусть $F_i=F_{i-1}[\beta]$, $\beta^q=b\in F_{i-1}$, $q$ простое.
\begin{itemize}
\item Если $q=p$, то минимальный многочлен $\beta$ над $F_{i-1}$ делит $X^p-b=(X-\beta)^p$, т.\,е.\ имеет единственный корень
$\beta$; поэтому $[F_i:F_{i-1}]_s=1$.
\item Пусть $q\ne p$. При $b=0$ шаг тривиален. Иначе, если $\zeta\in F$~--- первообразный корень степени $q$ из единицы,
то $X^q-b=\prod_{j=0}^{q-1}(X-\zeta^j\beta)$ с попарно различными корнями, все они лежат в $F_i$; значит, $F_i$~--- поле
разложения сепарабельного многочлена, $F_i/F_{i-1}$~--- расширение Галуа (\rf{3.10}). Для
$\sigma\in\Gal(F_i/F_{i-1})$ имеем $\sigma\beta=\zeta^{j}\beta$, и $\sigma\mapsto\sigma\beta/\beta$~--- гомоморфизм в группу
$\langle\zeta\rangle$ порядка $q$ (так как $\sigma\zeta=\zeta$), инъективный, поскольку $\beta$ порождает $F_i$ над $F_{i-1}$.
Поэтому $[F_i:F_{i-1}]_s=[F_i:F_{i-1}]=|\Gal(F_i/F_{i-1})|\in\{1,q\}$.
\end{itemize}
По мультипликативности $[F_m:F]_s$~--- произведение простых чисел, отличных от $p$; оно не делится на $p$.

С другой стороны, $F\subset F[\alpha]\subset F_m$. Так как $a\notin\wp(F)$, многочлен $X^p-X-a$ неприводим (\rf{5.32c}(b)) и
является минимальным многочленом $\alpha$, а его корни $\alpha+i$, $i\in\F_p$, попарно различны. Поэтому
$[F[\alpha]:F]_s=p$ и $[F_m:F]_s=[F_m:F[\alpha]]_s\cdot p$ делится на $p$~--- противоречие.

Для $F=\Fpb(t)$: многочлен $X^q-1$ при $q\ne p$ сепарабелен, так что в $\Fpb$ есть $q$ корней степени $q$ из единицы, и они
образуют циклическую группу порядка $q$ (\rf{5.8}); а $t\notin\wp(\Fpb(t))$ по \rf{5.32h}. (Другое доказательство, через
степени расширений Галуа внутри радикальной башни,~--- \rf{5.32i}.)
\end{proof}

\begin{center}
\begin{keepfig}
\centering
\begin{tikzpicture}[every node/.style={font=\small}]
\node (F) at (0,0) {$F$}; \node (Fa) at (-2.2,1.4) {$F[\alpha]$}; \node (Fm) at (0,2.8) {$F_m$};
\draw (F) -- node[left=2pt,cnew] {$p$} (Fa); \draw (Fa) -- (Fm);
\draw[cold] (F) -- node[right=2pt,align=left,font=\footnotesize] {радикальная башня:\\ шаги с $[\,\cdot\,]_s\in\{1,q\}$, $q\ne p$} (Fm);
\node[anchor=west,align=left,font=\footnotesize] at (5.2,1.4) {$[F_m:F]_s$ не делится на $p$,\\
но $[F_m:F]_s=[F_m:F[\alpha]]_s\cdot\textcolor{cnew}{p}$};
\end{tikzpicture}
\figcap{Доказательство \rf{3.29a}: число вложений вершины радикальной башни в $\Om$ не делится на $p$, а для поля,
содержащего корень неприводимого многочлена Артина--Шрайера, оно делилось бы на $p$.}
\end{keepfig}
\end{center}

Условие на корни из единицы в \rf{3.29a} существенно. Над конечным полем каждый многочлен разрешим в радикалах: поле
разложения $\F_{q^r}$ порождено образующей циклической группы $\F_{q^r}^\times$, т.\,е.\ корнем из единицы (\rf{5.35c}(a)). Например,
корни $X^p-X-1$ над $\F_p$~--- корни из единицы: из $\alpha^p=\alpha+1$ следует $\alpha^{p^i}=\alpha+i$, а произведение
корней $\prod_{i\in\F_p}(\alpha+i)$ равно $1$ (оно равно $(-1)^p f(0)=(-1)^{p+1}=1$ в характеристике $p$), так что
$\alpha^{1+p+\dots+p^{p-1}}=1$ (\rf{5.32g}). Над полем $\F_p(t)$, где корней из единицы почти нет, \rf{3.29a} применяют после
присоединения нужных корней из единицы; так доказывается, что $X^p-X-t$ не разрешим в радикалах и над $\F_p(t)$ (\rf{5.32j}).
Вообще, корни многочлена Артина--Шрайера выражаются в радикалах тогда и только тогда, когда они лежат в круговом
расширении $F[\zeta_N]$, $p\nmid N$ (см.\ после \rf{5.32j}).

\subsection*{Уравнения простой степени: критерии Галуа и Абеля}

\begin{aside}[3.30]
Абель (1828) сформулировал следующее утверждение: \emph{если три корня неприводимого уравнения простой степени связаны
так, что один из них рационально выражается через два других, то уравнение разрешимо в радикалах.} Силов (1902) считал
это утверждение неверным. Делинь (2021) доказал, что утверждение Абеля верно, если
понимать его для \emph{любых} трёх корней, и что верно и обратное. На языке теории Галуа всё сводится к группам перестановок: если $E$~---
множество из $p$ элементов ($p$ простое) и $G$~--- транзитивная группа перестановок $E$, то $G$ разрешима тогда и только
тогда, когда для любых трёх элементов $E$ один из них неподвижен относительно всех $g\in G$, оставляющих на месте два
других. Докажем это.
\end{aside}

Для $a,b\in E$ обозначим $G_a=\{g\in G\mid ga=a\}$ (стабилизатор точки $a$) и $G_{ab}=G_a\cap G_b=\{g\in G\mid ga=a,\ gb=b\}$. Через $\mathrm{AGL}_1(\F_p)$ обозначим группу аффинных
преобразований $x\mapsto ux+v$ ($u\in\F_p^\times$, $v\in\F_p$) множества $\F_p=\Z/p\Z$.

\begin{theorem}[3.30a: Галуа, Делинь]
Пусть $p$ простое, $E$~--- множество из $p$ элементов, $G$~--- транзитивная группа перестановок $E$. Следующие условия
равносильны:
\begin{enumerate}[label=\textup{(\roman*)}]
\item $G_{ab}=\{1\}$ для некоторых $a\ne b$;
\item $|G|<p^2$;
\item в $G$ есть нормальная подгруппа порядка $p$;
\item при подходящей нумерации $E=\F_p$ группа $G$ содержится в $\mathrm{AGL}_1(\F_p)$ и содержит все сдвиги $x\mapsto x+v$;
\item $G$ разрешима;
\item $G_{ab}=\{1\}$ для всех $a\ne b$, т.\,е.\ каждый элемент $g\ne1$ оставляет на месте не более одного элемента $E$;
\item \textup(условие Абеля\textup) для любых трёх различных $a,b,c\in E$ один из них неподвижен относительно всех
элементов $G$, оставляющих на месте два других.
\end{enumerate}
\end{theorem}

\begin{proof}
(i)$\Rightarrow$(ii). Орбита пары $(a,b)$ лежит в множестве пар различных элементов, поэтому
$|G|=|G\cdot(a,b)|\cdot|G_{ab}|\le p(p-1)<p^2$.

(ii)$\Rightarrow$(iii). Орбита транзитивной группы состоит из $p$ элементов, поэтому $p$ делит $|G|$, и по теореме
Коши в $G$ есть подгруппа $A$ порядка $p$. Если бы была другая подгруппа $B$ порядка $p$, то $A\cap B=\{1\}$, и
произведения $ab$ ($a\in A$, $b\in B$) были бы попарно различны ($ab=a'b'\Rightarrow a'^{-1}a=b'b^{-1}\in A\cap B$), т.\,е.\ в
$G$ было бы не меньше $p^2$ элементов. Итак, $A$~--- единственная подгруппа порядка $p$; так как $gAg^{-1}$~--- тоже
подгруппа порядка $p$, она нормальна.

(iii)$\Rightarrow$(iv). Пусть $A=\langle h\rangle$, $|A|=p$. Длины циклов перестановки $h$ делят её порядок $p$ и не все равны
$1$, а сумма длин равна $p$; значит, $h$~--- цикл длины $p$. Занумеруем $E$ элементами $\F_p$ вдоль этого цикла: $h(n)=n+1$;
тогда $A$~--- группа всех сдвигов. Для $g\in G$ элемент $ghg^{-1}\in A$ отличен от $1$, т.\,е.\ $ghg^{-1}=h^u$, $u\in\F_p^\times$.
Значит, $g(n+1)=g(h(n))=h^u(g(n))=g(n)+u$, и по индукции $g(n)=un+v$, где $v=g(0)$.

(iv)$\Rightarrow$(v). Отображение $(x\mapsto ux+v)\mapsto u$~--- гомоморфизм $\mathrm{AGL}_1(\F_p)\to\F_p^\times$ (композиция
$x\mapsto u(u'x+v')+v$ имеет коэффициент $uu'$) с ядром $T$ из сдвигов. Группа $T\cong\Z/p\Z$ абелева, факторгруппа
вкладывается в абелеву группу $\F_p^\times$; поэтому $\mathrm{AGL}_1(\F_p)$ разрешима, а с ней и её подгруппа $G$.

(iv)$\Rightarrow$(vi). Если $ux+v=x$ для двух различных $x$, то, вычитая, получаем $u=1$, и тогда $v=0$.

(vi)$\Rightarrow$(i) и (vi)$\Rightarrow$(vii) очевидны: элементы $G_{ab}=\{1\}$ неподвижны на всём $E$.

(vii)$\Rightarrow$(i). Выберем пару $a\ne b$, для которой $|G_{ab}|$ \emph{наименьший}, и пусть $c$~--- любой третий элемент.
По (vii) верно одно из трёх: $G_{ab}$ оставляет $c$ на месте; или $G_{ac}\subset G_b$, т.\,е.\ $G_{ac}\subset G_{ab}$, но по выбору
пары $|G_{ac}|\ge|G_{ab}|$, так что $G_{ac}=G_{ab}$, и $G_{ab}$ снова оставляет $c$ на месте; или $G_{bc}\subset G_a$, и
так же $G_{bc}=G_{ab}$ оставляет $c$ на месте. Итак, $G_{ab}$ неподвижна на всём $E$, т.\,е.\ $G_{ab}=\{1\}$.

(v)$\Rightarrow$(iii). Пусть $\{1\}=H_0\subsetneq H_1\subsetneq\dots\subsetneq H_n=G$, $H_{i-1}$ нормальна в $H_i$, и $H_i/H_{i-1}$
абелевы (индексы здесь~--- номера членов цепочки, а не точки $E$). \emph{Нетривиальная нормальная подгруппа $N$
транзитивной группы $K$ транзитивна}: $k(Nx)=N(kx)$ для $k\in K$, так что $K$ транзитивно переставляет $N$-орбиты, они все
одного размера, и этот размер делит $p$; размер $1$ означал бы, что $N$ действует тождественно, т.\,е.\ $N=\{1\}$. По
убывающей индукции все $H_i$, $i\ge1$, транзитивны. Группа $H_1$ абелева и транзитивна, поэтому $|H_1|=p$: стабилизаторы
точек в ней сопряжены, т.\,е.\ (по абелевости) совпадают, так что стабилизатор точки неподвижен на всём $E$ и тривиален, и
$|H_1|$ равен длине орбиты. Положим $A=H_1$. Она нормальна в $H_2$; по уже доказанному (iii)$\Rightarrow$(iv) группа $H_2$
лежит в $\mathrm{AGL}_1(\F_p)$, так что $|H_2|\le p(p-1)<p^2$, и, как в (ii)$\Rightarrow$(iii), $A$~--- единственная подгруппа
порядка $p$ в $H_2$. Для $g\in H_3$ имеем $gH_2g^{-1}=H_2$, поэтому $gAg^{-1}$~--- подгруппа порядка $p$ в $H_2$, т.\,е.\
$gAg^{-1}=A$: $A$ нормальна в $H_3$. Повторяя рассуждение, получаем, что $A$ нормальна в $H_n=G$.
\end{proof}

\begin{corollary}[3.30b]
Пусть $\operatorname{char}F=0$, $f\in F[X]$~--- неприводимый многочлен простой степени $p$, $E$~--- его поле разложения.
Следующие условия равносильны:
\begin{enumerate}[label=\textup{(\alph*)}]
\item $f$ разрешим в радикалах;
\item $E=F[\alpha,\beta]$ для любых двух различных корней $\alpha,\beta$ многочлена $f$ \textup(теорема Галуа, \rf{5.35}\textup);
\item $E=F[\alpha,\beta]$ для некоторых двух различных корней $\alpha,\beta$;
\item \textup(критерий Абеля\textup) из любых трёх различных корней $f$ один лежит в поле, порождённом над $F$ двумя
другими, т.\,е.\ рационально выражается через них с коэффициентами из $F$.
\end{enumerate}
\end{corollary}

\begin{proof}
Группа $G=G_f$ действует на множестве из $p$ корней транзитивно (\rf{3.27d}). Для корней $\alpha,\beta$ элемент $G$
тождествен на $F[\alpha,\beta]$ тогда и только тогда, когда он оставляет на месте $\alpha$ и $\beta$; поэтому
$\Gal(E/F[\alpha,\beta])=G_{\alpha\beta}$, и по основной теореме (\rf{3.17}) $F[\alpha,\beta]=E^{G_{\alpha\beta}}$. Значит,
корень $\gamma$ лежит в $F[\alpha,\beta]$ тогда и только тогда, когда $G_{\alpha\beta}$ оставляет $\gamma$ на месте, а
$E=F[\alpha,\beta]$ тогда и только тогда, когда $G_{\alpha\beta}=\{1\}$. Поэтому (b), (c), (d)~--- это условия (vi), (i), (vii)
теоремы \rf{3.30a}, а (a) по теореме \rf{3.28} равносильно (v).
\end{proof}

Например, у $X^5-2$ любые два корня порождают поле разложения (\rf{5.35}). А у многочлена $X^5-6X+3$ с группой $S_5$
(\rf{5.35e}) стабилизатор двух корней~--- группа $S_3$ перестановок трёх остальных; она не оставляет на месте ни одного из
них, так что поле $\Q[\alpha,\beta]$ степени $120/6=20$ не содержит ни одного из трёх других корней.

\begin{corollary}[3.30c: Кронекер]
Пусть $F\subset\R$, и $f\in F[X]$~--- неприводимый многочлен нечётной простой степени $p$, разрешимый в радикалах. Тогда у
$f$ либо ровно один вещественный корень, либо все корни вещественны.
\end{corollary}

\begin{proof}
Вещественный корень есть, так как степень нечётна. Если $\alpha\ne\beta$~--- два вещественных корня, то по \rf{3.30b}(b) поле
разложения $F[\alpha,\beta]\subset\CC$ лежит в $\R$, т.\,е.\ все корни вещественны.
\end{proof}

Так, у $X^5-6X+3$ ровно три вещественных корня (\rf{5.35e}), и по \rf{3.30c} он не разрешим в радикалах~--- без вычисления
группы Галуа.

\begin{example}[3.30d: почему нужны \emph{любые} три корня]
Пусть $E=\F_2^3\smallsetminus\{0\}$~--- семь точек \emph{плоскости Фано}; её <<прямые>>~--- тройки $\{a,b,a+b\}$ (их семь).
Группа $G=\mathrm{GL}_3(\F_2)$ действует на $E$ транзитивно (любой ненулевой вектор дополняется до базиса), и
$|G|=(8-1)(8-2)(8-4)=168$ (первый столбец обратимой матрицы~--- любой ненулевой вектор, второй~--- не лежащий на
порождённой первым прямой, третий~--- не лежащий в плоскости первых двух). Так как $168\ge7^2$, по \rf{3.30a} группа $G$
неразрешима. Тем не менее условие из (vii) выполнено для каждой \emph{прямой} $\{a,b,a+b\}$: линейное отображение,
оставляющее на месте $a$ и $b$, оставляет на месте и $a+b$. Для тройки $\{a,b,c\}$, не лежащей на прямой (базис), оно
нарушено: линейные отображения, задаваемые на базисе правилами $c\mapsto c+a$ (остальные на месте), $b\mapsto b+a$ и
$a\mapsto a+b$, обратимы, оставляют на месте по два элемента тройки и сдвигают третий. Поэтому если группа Галуа
неприводимого многочлена седьмой степени действует на его корнях как $G$ на точках плоскости Фано, то найдутся три корня,
один из которых рационально выражается через два других (отвечающие прямой), но многочлен не разрешим в радикалах.
Таков, например, $X^7-7X+3$ над $\Q$: он неприводим (по модулю $2$ это неприводимый многочлен $X^7+X+1$), имеет ровно
три вещественных корня и потому по \rf{3.30c} не разрешим в радикалах, а его группа Галуа (это проверяется
вычислением) имеет порядок $168$ и действует на корнях как $\mathrm{GL}_3(\F_2)$ на точках плоскости Фано.

Так же ведут себя группы $\mathrm{GL}_k(\F_2)$ на $2^k-1$ ненулевых векторах пространства $\F_2^k$ при $k\ge3$, когда
число $2^k-1$ простое: $|\mathrm{GL}_k(\F_2)|\ge(2^k-1)(2^k-2)(2^k-4)>(2^k-1)^2$, тройки $\{a,b,a+b\}$ удовлетворяют
условию, а тройки линейно независимых векторов~--- нет. (При $k=2$ группа $\mathrm{GL}_2(\F_2)\cong S_3$ разрешима, и
любые три ненулевых вектора $\F_2^2$ образуют прямую.)
\end{example}

\begin{center}
\begin{keepfig}
\centering
\begin{tikzpicture}[scale=1.25,every node/.style={font=\scriptsize}]
\coordinate (A) at (90:2); \coordinate (B) at (210:2); \coordinate (C) at (330:2);
\coordinate (AB) at ($(A)!0.5!(B)$); \coordinate (AC) at ($(A)!0.5!(C)$); \coordinate (BC) at ($(B)!0.5!(C)$);
\coordinate (O) at (0,0);
\draw[gray] (A) -- (C); \draw[gray] (B) -- (C);
\draw[gray] (A) -- (BC); \draw[gray] (B) -- (AC); \draw[gray] (C) -- (AB);
\draw[gray] (O) circle (1);
\draw[cnew,very thick] (A) -- (B);
\foreach \P in {A,B,C,AB,AC,BC,O}{\fill (\P) circle (1.6pt);}
\fill[cnew] (A) circle (2pt); \fill[cnew] (B) circle (2pt); \fill[cnew] (AB) circle (2pt);
\draw[cold,thick] (A) circle (4pt); \draw[cold,thick] (B) circle (4pt); \draw[cold,thick] (C) circle (4pt);
\node[above=3pt] at (A) {$100$}; \node[below left=1pt] at (B) {$010$}; \node[below right=1pt] at (C) {$001$};
\node[left=3pt] at (AB) {$110$}; \node[right=3pt] at (AC) {$101$}; \node[below=3pt] at (BC) {$011$};
\node[right=5pt] at (O) {$111$};
\end{tikzpicture}
\figcap{Плоскость Фано: семь точек $\F_2^3\smallsetminus\{0\}$ и семь прямых (шесть отрезков и окружность). Для тройки на
прямой (красная) условие Абеля выполнено, для тройки не на прямой (обведённые синим вершины) нарушено.}
\end{keepfig}
\end{center}
